\scanpage{148}
\mspage{V.20}
% archive sheet number: 147

Dans ce dernier cas, \deleted{$\unsure{\bar V\longrightarrow}$} l'inclusion $V_{\unsure{\bar\eta|\bar y}}\longrightarrow\bar V$ s'apparente à un isomorphisme du pt de vue cohomologique (coeff. des $\Lambda_{\unsure{\nu|v}}$), et on peut dire que les cinq schémas envisagés \deleted{\unsure{se sont}} \deleted{\unsure{\&}} se réduisent à quatre, $V_s,V,V_{\unsure{\hat\eta|\bar\eta|\eta}},V_{\bar\eta}$ (on est débarrassé du seul $\bar V$ de ces schémas qui n'est pas noethérien !).

4.5. L'inclusion $V_{\bar s}\lhook\joinrel\longrightarrow\bar V$ définit de $\unsure{=|\simeq|\bar m}$ un homomorphisme
\[
\tag*{(4.5.1)}
\unsure{R|\mathbb R}(\Gamma,I)(\bar V,\Lambda_\nu)\longrightarrow\mathbb R(\Gamma,I)(V_s,\Lambda_\nu)
\]
qui s'insère dans un triangle exact

\deleted{\unsure{(4.5.2)}} (B)
\diagram{\begin{tikzcd}[ampersand replacement=\&,column sep=large,row sep=large]
\& \mathbb R(\Gamma,I)(V_{\unsure{\bar s|s}},\Lambda_\nu)\simeq\mathbb R\Gamma(V_s,\Lambda_\nu)^{\mathrm{triv}}\arrow[dl,"{\unsure{\partial|\delta|)}}"'] \& \\
\mathbb R(\Gamma,I)(\bar V,(\Lambda_\nu)_{V_{\bar\eta},V})\arrow[rr] \& \& \mathbb R(\Gamma,I)(\bar V,\Lambda_\nu)\arrow[ul,"\bar i^*"]
\end{tikzcd}}.

L'application du foncteur oubli à ce diagramme donne

(B')
\diagram{\begin{tikzcd}[ampersand replacement=\&,column sep=large,row sep=large]
\& \mathbb R\Gamma(V_s,\Lambda_\nu)\arrow[dl,"\partial"'] \& \\
\mathbb R\Gamma(\bar V,(\Lambda_\nu)_{V_{\bar\eta},V})\arrow[rr] \& \& \mathbb R\Gamma(\bar V,\Lambda_\nu)\arrow[ul,"\bar i^*"]
\end{tikzcd}}.

\deleted{Sous la condition (4.5.2)}
